% Hafta 4 -- alıştırma gövdesi (ders notu ve ayrı alıştırma kağıdı tarafından \input edilir)
\section{Çözümlü Alıştırmalar}
Her soruyu önce kendiniz deneyin, sonra çözümü okuyun.

\begin{alistirma}
\(\left(1\dfrac{1}{2}+\dfrac{3}{4}\right)\div 1\dfrac{1}{8}\) işleminin sonucu kaçtır?
\end{alistirma}
\begin{cozum}
Bileşik kesre çevirelim: \(1\frac{1}{2}=\frac{3}{2}\), \(1\frac{1}{8}=\frac{9}{8}\). \ Parantez: \(\dfrac{3}{2}+\dfrac{3}{4}=\dfrac{6}{4}+\dfrac{3}{4}=\dfrac{9}{4}\). \quad Bölme: \(\dfrac{9}{4}\div\dfrac{9}{8}=\dfrac{9}{4}\cdot\dfrac{8}{9}=\dfrac{8}{4}=2\).
\end{cozum}

\begin{alistirma}
\(3-\cfrac{2}{1+\cfrac{1}{2}}\) işleminin sonucu kaçtır?
\end{alistirma}
\begin{cozum}
En alttan başlanır: \(1+\dfrac{1}{2}=\dfrac{3}{2}\) \(\ \rightarrow\ \) \(\dfrac{2}{3/2}=2\cdot\dfrac{2}{3}=\dfrac{4}{3}\) \(\ \rightarrow\ \) \(3-\dfrac{4}{3}=\dfrac{9}{3}-\dfrac{4}{3}=\dfrac{5}{3}\).
\end{cozum}

\begin{alistirma}
\(1{,}25\cdot 0{,}08+0{,}9\div 0{,}03\) işleminin sonucu kaçtır?
\end{alistirma}
\begin{cozum}
Önce çarpma ve bölme, sonra toplama. Çarpma: \(125\cdot 8=1000\); ondalık basamak sayısı \(2+2=4\) \(\Rightarrow\) \(0{,}1000=0{,}1\).\\
Bölme: pay ve payda \(100\) ile genişletilir: \(\dfrac{0{,}9}{0{,}03}=\dfrac{90}{3}=30\). \quad Toplam: \(0{,}1+30=30{,}1\).
\end{cozum}

\begin{alistirma}
\(1{,}\ab{6}+0{,}8\ab{3}\) toplamı kaçtır?
\end{alistirma}
\begin{cozum}
Devirli sayılar önce kesre çevrilir:\\
\(1{,}\ab{6}=\dfrac{16-1}{9}=\dfrac{15}{9}=\dfrac{5}{3}\) \ (bir basamak devrediyor: bir \(9\)); \quad
\(0{,}8\ab{3}=\dfrac{83-8}{90}=\dfrac{75}{90}=\dfrac{5}{6}\) \ (bir \(9\), bir \(0\)).\\
Toplam: \(\dfrac{5}{3}+\dfrac{5}{6}=\dfrac{10}{6}+\dfrac{5}{6}=\dfrac{15}{6}=\dfrac{5}{2}=2{,}5\).
\end{cozum}

\begin{alistirma}
\(a=-\dfrac{3}{5}\), \(b=-0{,}\ab{6}\), \(c=-\dfrac{7}{12}\) sayılarını küçükten büyüğe sıralayınız.
\end{alistirma}
\begin{cozum}
\(0{,}\ab{6}=\dfrac{6}{9}=\dfrac{2}{3}\). Önce pozitifler sıralanır; \(\text{EKOK}(5,3,12)=60\): \ \(\dfrac{3}{5}=\dfrac{36}{60}\), \(\dfrac{2}{3}=\dfrac{40}{60}\), \(\dfrac{7}{12}=\dfrac{35}{60}\) \(\Rightarrow\) \(\dfrac{7}{12}<\dfrac{3}{5}<\dfrac{2}{3}\).\\
Negatiflerde sıra ters döner: \(-\dfrac{2}{3}<-\dfrac{3}{5}<-\dfrac{7}{12}\), yani \(b<a<c\).
\end{cozum}

\begin{alistirma}
\(\dfrac{3}{8}\) ile \(\dfrac{1}{2}\) arasında olup \(\dfrac{k}{24}\ (k\in\mathbb{Z})\) biçiminde yazılabilen kaç rasyonel sayı vardır?
\end{alistirma}
\begin{cozum}
Paydalar \(24\)'e genişletilir: \(\dfrac{3}{8}=\dfrac{9}{24}\), \ \(\dfrac{1}{2}=\dfrac{12}{24}\).\\
\(\dfrac{9}{24}<\dfrac{k}{24}<\dfrac{12}{24}\iff 9<k<12\) \(\Rightarrow\) \(k\in\{10,11\}\): \ \(2\) tane (\(\frac{10}{24}=\frac{5}{12}\) ve \(\frac{11}{24}\)).
\end{cozum}

\begin{alistirma}
\(720\) TL'nin \(\frac{5}{12}\)'si harcanıyor, kalan paranın \(\frac{3}{7}\)'si bağışlanıyor. Geriye kaç TL kalır?
\end{alistirma}
\begin{cozum}
Harcanan: \(720\cdot\dfrac{5}{12}=300\) TL \(\Rightarrow\) kalan \(720-300=420\) TL.\\
Bağış (``kalanın'' \(\frac{3}{7}\)'si): \(420\cdot\dfrac{3}{7}=180\) TL \(\Rightarrow\) geriye \(420-180=240\) TL. \ Kısa yol: \(720\cdot\dfrac{7}{12}\cdot\dfrac{4}{7}=240\) TL.
\end{cozum}

\begin{alistirma}
\(\dfrac{1}{2\cdot 3}+\dfrac{1}{3\cdot 4}+\dots+\dfrac{1}{11\cdot 12}\) toplamı kaçtır?
\end{alistirma}
\begin{cozum}
\(\dfrac{1}{n(n+1)}=\dfrac{1}{n}-\dfrac{1}{n+1}\) olduğundan: \ \(\left(\frac{1}{2}-\frac{1}{3}\right)+\left(\frac{1}{3}-\frac{1}{4}\right)+\dots+\left(\frac{1}{11}-\frac{1}{12}\right)\).\\
Ara terimler birbirini götürür: \(\dfrac{1}{2}-\dfrac{1}{12}=\dfrac{6}{12}-\dfrac{1}{12}=\dfrac{5}{12}\).
\end{cozum}

\kendinizoncesi
\section{Alıştırmalar (Kendiniz Çözün)}

\begin{enumerate}[label=\textbf{\arabic*.},leftmargin=*,itemsep=3pt]
\item \(6\tfrac{2}{3}\) sayısını bileşik kesir, \(\tfrac{47}{6}\) kesrini tam sayılı kesir olarak yazınız.
\item \(\tfrac{126}{168}\) kesrinin en sade biçimi nedir?
\item \(0{,}64\) ve \(2{,}05\) sayılarını en sade kesir olarak yazınız.
\item \(\tfrac{3}{8}+\tfrac{5}{12}-\tfrac{1}{6}\) işleminin sonucu kaçtır?
\item \(\tfrac{2}{3}\cdot\tfrac{9}{10}\div\tfrac{3}{5}\) işleminin sonucu kaçtır?
\item \(\tfrac{0{,}056}{0{,}0014}\) işleminin sonucu kaçtır?
\item \(0{,}00072\) ve \(58\,000\,000\) sayılarını bilimsel gösterimle yazınız.
\item \(0{,}\ab{36}\) ve \(0{,}4\ab{5}\) sayılarını en sade kesir olarak yazınız.
\item \(0{,}\ab{3}+0{,}\ab{6}\cdot 0{,}75\) işleminin sonucunu en sade kesir olarak yazınız.
\item \(a=\tfrac{7}{9}\), \(b=\tfrac{9}{11}\), \(c=\tfrac{5}{7}\) sayılarını küçükten büyüğe sıralayınız.
\item \(\tfrac{5}{6}\) ile \(\tfrac{7}{8}\) sayılarının tam ortasındaki (aritmetik ortalaması olan) sayı kaçtır?
\item \(840\) litrelik bir deponun \(\frac{3}{7}\)'si doludur. Depoya önce deponun \textbf{tamamının} \(\frac{1}{4}\)'ü kadar su ekleniyor, sonra depodaki suyun \(\frac{1}{3}\)'ü kullanılıyor. Depoda kaç litre su kalır? Bu miktar deponun kaçta kaçıdır?
\end{enumerate}

\subsection*{Cevap Anahtarı}
{\small
\begin{multicols}{2}
\begin{enumerate}[label=\textbf{\arabic*.},leftmargin=*,itemsep=1pt]
\item \(\frac{6\cdot 3+2}{3}=\frac{20}{3}\); \ \(47=6\cdot 7+5\Rightarrow 7\frac{5}{6}\).
\item \(\text{EBOB}(126,168)=42\): \(\frac{3}{4}\).
\item \(\frac{64}{100}=\frac{16}{25}\); \ \(\frac{205}{100}=\frac{41}{20}\).
\item \(\text{EKOK}=24\): \(\frac{9+10-4}{24}=\frac{15}{24}=\frac{5}{8}\).
\item Soldan sağa: \(\frac{3}{5}\div\frac{3}{5}=1\).
\item \(10000\) ile genişlet: \(\frac{560}{14}=40\).
\item \(7{,}2\cdot 10^{-4}\); \ \(5{,}8\cdot 10^{7}\).
\item \(\frac{36}{99}=\frac{4}{11}\); \ \(\frac{45-4}{90}=\frac{41}{90}\).
\item \(\frac{1}{3}+\frac{2}{3}\cdot\frac{3}{4}=\frac{1}{3}+\frac{1}{2}=\frac{5}{6}\).
\item \(1\)'e farklar \(\frac{2}{9},\frac{2}{11},\frac{2}{7}\): \(c<a<b\).
\item \(\frac{1}{2}\left(\frac{20}{24}+\frac{21}{24}\right)=\frac{41}{48}\).
\item \(360+210=570\); \(570\cdot\frac{2}{3}=380\) L; \(\frac{380}{840}=\frac{19}{42}\).
\end{enumerate}
\end{multicols}}
